%% =====================================================================
%%  Mode d'emploi — le code couleur des fiches
%% =====================================================================
\thispagestyle{otFront}
\pdfbookmark[0]{Mode d'emploi}{bm-mode-emploi}
\otPageTitre{Mode d'emploi}{Comment lire une fiche}

{\fontsize{10.5}{14}\selectfont
Chaque fiche condense un chapitre du programme en \textbf{trois à quatre pages} : tout ce qu'il
faut savoir et savoir faire le jour du bac, rien de plus, rien de moins. Les notions sont rangées
par couleur, pour retrouver d'un coup d'œil ce que l'on cherche.\par}
\vspace{4mm}

\noindent
\begin{minipage}[t]{0.485\linewidth}
\begin{definition}
Le vocabulaire exact du cours, à connaître mot pour mot.
\end{definition}
\begin{theoreme}
Les grands résultats du programme, avec toutes leurs hypothèses.
\end{theoreme}
\begin{propriete}
Les résultats d'usage courant, à savoir appliquer sans hésiter.
\end{propriete}
\formule{\text{Les formules à savoir par cœur}}
\begin{demonstration}
Les preuves exigibles par le programme, résumées en quelques lignes.
\end{demonstration}
\end{minipage}\hfill
\begin{minipage}[t]{0.485\linewidth}
\begin{methode}
La démarche type, étape par étape, pour réussir les exercices classiques.
\end{methode}
\begin{attention}
Les pièges et les erreurs les plus fréquentes dans les copies.
\end{attention}
\begin{exemple}
Une application directe, entièrement calculée.
\end{exemple}
\begin{redaction}
Le modèle de rédaction attendu sur la copie, à reproduire.
\end{redaction}
\begin{tcolorbox}[otBase,colback=otNavy,colframe=otNavy,coltext=white]
\noindent{\otSemi\fontsize{6.8}{8}\selectfont\color{otGreen}\otCaps[11]{Check-list}}\par\vspace{1.1mm}
En fin de fiche, pour vérifier que tout est acquis.
\end{tcolorbox}
\end{minipage}

\vspace{6mm}
\otIntertitre{Réviser efficacement avec les fiches}
\vspace{1mm}
{\fontsize{10.2}{13.6}\selectfont
\newcommand{\otConseilF}[2]{\par\noindent\makebox[8mm][l]{\tikz[baseline=-0.65ex]{\fill[otNavy] (0,0) circle (2.2mm);
  \node[text=white,font=\otSemi\fontsize{7.5}{8}\selectfont] at (0,0) {#1};}}%
  \parbox[t]{\dimexpr\linewidth-8mm\relax}{#2}\par\vspace{2.6mm}}
\otConseilF{1}{\textbf{Lisez la fiche avant les exercices} : elle rappelle l'essentiel du cours en dix minutes.}
\otConseilF{2}{\textbf{Cachez les encadrés verts} et retrouvez la méthode seul, puis vérifiez.}
\otConseilF{3}{\textbf{Cochez la check-list} : ce qui n'est pas coché est à retravailler en priorité.}
\otConseilF{4}{\textbf{Entraînez-vous} avec les exercices corrigés du livre OptiTAB du même chapitre :
la fiche et le chapitre portent le même numéro.}}
\clearpage
